\ExamPart{Câu tự luận}{3}{}

\NQuestion Giải bất phương trình
\[
\left(x^2-5x+6\right)\left(x^2-1\right)\le 0.
\]
\AnswerLine{$[-1;1]\cup[2;3]$}
\SolutionBlock{Ta có
\[
\left(x^2-5x+6\right)\left(x^2-1\right)=(x-2)(x-3)(x-1)(x+1).
\]
Các nghiệm là $-1,1,2,3$. Lập bảng xét dấu theo thứ tự các nghiệm, vì hệ số bậc cao nhất dương nên dấu của biểu thức lần lượt là $+,-,+,-,+$. Do đó bất phương trình đúng khi
\[
x\in[-1;1]\cup[2;3].
\]}

\NQuestion Tìm tất cả các giá trị của $m$ để phương trình
\[
x^2-2(m+1)x+m^2+3=0
\]
có hai nghiệm phân biệt $x_1,x_2$ thỏa mãn $x_1+x_2>2$.
\AnswerLine{$m>1$}
\SolutionBlock{Theo hệ thức Viète,
\[
x_1+x_2=2(m+1).
\]
Điều kiện $x_1+x_2>2$ cho $2(m+1)>2\Rightarrow m>0$. Để phương trình có hai nghiệm phân biệt, cần và đủ $\Delta>0$. Ta có
\[
\Delta=[-2(m+1)]^2-4(m^2+3)=4(m+1)^2-4m^2-12=8m-8.
\]
Vậy $\Delta>0\Leftrightarrow m>1$. Kết hợp với $m>0$, suy ra $m>1$.}

\NQuestion Trong mặt phẳng tọa độ $Oxy$, cho ba điểm $A(1;2),\ B(5;4),\ C(3;-2)$.
\begin{SubQuestions}
  \item Tính tọa độ các vectơ $\overrightarrow{AB}$ và $\overrightarrow{AC}$.
  \item Chứng minh rằng tam giác $ABC$ vuông tại $A$.
  \item Tính diện tích tam giác $ABC$.
\end{SubQuestions}
\AnswerLine{$\overrightarrow{AB}=(4;2),\ \overrightarrow{AC}=(2;-4),\ S_{ABC}=10$}
\SolutionBlock{\begin{SubQuestions}
\item Ta có
\[
\overrightarrow{AB}=(5-1;4-2)=(4;2),\qquad \overrightarrow{AC}=(3-1;-2-2)=(2;-4).
\]
\item Tính tích vô hướng:
\[
\overrightarrow{AB}\cdot\overrightarrow{AC}=4\cdot2+2\cdot(-4)=8-8=0.
\]
Vì tích vô hướng bằng $0$ nên $\overrightarrow{AB}\perp\overrightarrow{AC}$. Vậy tam giác $ABC$ vuông tại $A$.
\item Diện tích tam giác là
\[
S_{ABC}=\dfrac12\left|\begin{vmatrix}4&2\\2&-4\end{vmatrix}\right|=\dfrac12|4\cdot(-4)-2\cdot2|=\dfrac12\cdot20=10.
\]
\end{SubQuestions}}

\NQuestion Trong mặt phẳng tọa độ $Oxy$, cho điểm $A(1;2)$ và đường thẳng
\[
d:2x+5y-9=0.
\]
\begin{SubQuestions}
  \item Viết phương trình đường thẳng $\Delta$ đi qua $A$ và vuông góc với $d$.
  \item Gọi $H$ là giao điểm của $\Delta$ và $d$. Tính tọa độ điểm $H$ và khoảng cách từ $A$ đến đường thẳng $d$.
\end{SubQuestions}
\AnswerLine{$\Delta:5x-2y-1=0,\ H\left(\dfrac{23}{29};\dfrac{43}{29}\right),\ d(A,d)=\dfrac{3}{\sqrt{29}}$}
\SolutionBlock{\begin{SubQuestions}
\item Đường thẳng $d:2x+5y-9=0$ có vectơ pháp tuyến là $(2;5)$. Vì $\Delta\perp d$ nên có thể lấy vectơ pháp tuyến của $\Delta$ là $(5;-2)$. Vậy $\Delta$ có dạng $5x-2y+c=0$. Thay $A(1;2)$ vào, ta được $5\cdot1-2\cdot2+c=0\Rightarrow c=-1$. Suy ra
\[
\Delta:5x-2y-1=0.
\]
\item Tọa độ $H$ là nghiệm của hệ
\[
\begin{cases}
2x+5y-9=0,\\
5x-2y-1=0.
\end{cases}
\]
Nhân phương trình đầu với $2$ và phương trình sau với $5$ rồi cộng lại:
\[
\begin{cases}
4x+10y=18,\\
25x-10y=5.
\end{cases}
\Rightarrow 29x=23\Rightarrow x=\dfrac{23}{29}.
\]
Thế vào $2x+5y=9$ được $y=\dfrac{43}{29}$. Do đó
\[
H\left(\dfrac{23}{29};\dfrac{43}{29}\right).
\]
Khoảng cách từ $A$ đến $d$ là
\[
d(A,d)=\dfrac{|2\cdot1+5\cdot2-9|}{\sqrt{2^2+5^2}}=\dfrac{3}{\sqrt{29}}.
\]
\end{SubQuestions}}
